% Sección: Parte 2 – ¿Qué son las PINNs?

% Frame 1: Definición y motivación
\begin{frame}{¿Qué son las PINNs?}
  \begin{itemize}
    \item Physics-Informed Neural Networks (PINNs): redes neuronales que incorporan términos de las ecuaciones físicas en la función de pérdida.
    \item Combina datos (si disponibles) con conocimiento de leyes físicas (Ecuaciones diferenciales, condiciones de contorno).
    \item Objetivo: aproximar soluciones a ecuaciones diferenciales sin necesidad de generadores de malla tradicionales.
  \end{itemize}
\end{frame}

% Frame 2: ¿Por qué usar PINNs? ¿Cómo ayudan?
\begin{frame}{¿Por qué y cómo nos ayudan las PINNs?}
  \begin{itemize}
    \item Sin mallas: evitan el mallado costoso de dominio complejo.
    \item Generalización: una vez entrenada, la red evalúa la solución en cualquier punto sin recalcular todo el dominio.
    \item Integración de datos experimentales: robustez y ajuste basado en observaciones.
    \item Flexibilidad: adaptación a condiciones de frontera variadas y parámetros cambiantes.
  \end{itemize}
\end{frame}

% Frame 3: Línea de tiempo de las PINNs
\begin{frame}{Línea de tiempo de las PINNs}
  \begin{itemize}
    \item 2017: Raissi, Perdikaris & Karniadakis introducen el concepto en «Physics-informed neural networks: A deep learning framework for solving forward and inverse problems involving nonlinear partial differential equations.»
    \item 2019: Extensiones a problemas de mecánica de fluidos y transferencia de calor.
    \item 2020–2021: Desarrollo de librerías especializadas (DeepXDE, Modulus).
    \item 2022–2024: Aplicaciones en mecánica cuántica, biología computacional y geociencias.
  \end{itemize}
\end{frame}

% Frame 4: Ventajas y desventajas
\begin{frame}{Ventajas y desventajas de las PINNs}
  \begin{columns}
    \column{0.48\textwidth}
    \textbf{Ventajas:}
    \begin{itemize}
      \item Reducción de costos de mallado.
      \item Capacidad de incorporar datos ruidosos.
      \item Evaluación rápida tras entrenamiento.
      \item Enfoque unificado para problemas directos e inversos.
    \end{itemize}
    \column{0.48\textwidth}
    \textbf{Desventajas:}
    \begin{itemize}
      \item Entrenamiento costoso (optimización de funciones de pérdida complejas).
      \item Requerimiento de hiperparámetros finos (arquitectura, parámetros de la pérdida).
      \item Dificultad para problemas muy rígidos o con condiciones de contorno complicadas.
      \item Riesgo de sobreajuste si hay pocos datos.
    \end{itemize}
  \end{columns}
\end{frame}

